\chapter{Sequence Models}

\section{Recurrent Neural Networks}

Many real-world data types are sequential, including time series, text, and speech. Sequence models aim to capture dependencies across time.

\medskip

\noindent
\textbf{Recurrent neural networks (RNNs).}  
An RNN processes a sequence $(x_1, x_2, \dots, x_T)$ by maintaining a hidden state $h_t$ that evolves over time:
\[
h_t = \phi(W_h h_{t-1} + W_x x_t + b),
\]
\[
y_t = \psi(W_y h_t + c),
\]
where $\phi$ and $\psi$ are nonlinear functions.

\medskip

\noindent
\textbf{Interpretation.}
\begin{itemize}
    \item $h_t$ summarizes past information
    \item The same parameters are reused across time steps
\end{itemize}

\medskip

\noindent
\textbf{Parameter sharing.}  
Unlike feedforward networks, RNNs share weights across time, making them suitable for variable-length sequences.

\medskip

\noindent
\textbf{Unrolling in time.}  
An RNN can be viewed as a deep network with shared parameters:
\[
h_1 \rightarrow h_2 \rightarrow \cdots \rightarrow h_T.
\]

\medskip

\noindent
\textbf{Insight.}  
RNNs model temporal dependencies by recursively updating a hidden state.

\section{Gated Architectures (LSTM, GRU)}

Standard RNNs struggle to capture long-term dependencies. Gated architectures address this limitation.

\subsection{Long Short-Term Memory (LSTM)}

\medskip

\noindent
\textbf{Core idea.}  
Introduce a memory cell that can preserve information over long time periods.

\medskip

\noindent
\textbf{Gates.}
\begin{itemize}
    \item Input gate: controls new information
    \item Forget gate: controls what to discard
    \item Output gate: controls what to expose
\end{itemize}

\medskip

\noindent
\textbf{Structure (conceptual).}
\begin{itemize}
    \item Cell state evolves with controlled updates
    \item Gates regulate information flow
\end{itemize}

\medskip

\noindent
\textbf{Effect.}
\begin{itemize}
    \item Enables learning of long-term dependencies
    \item Mitigates vanishing gradient problem
\end{itemize}

\subsection{Gated Recurrent Unit (GRU)}

\medskip

\noindent
\textbf{Simplified architecture.}
\begin{itemize}
    \item Combines input and forget mechanisms
    \item Fewer parameters than LSTM
\end{itemize}

\medskip

\noindent
\textbf{Update mechanism.}
\begin{itemize}
    \item Update gate controls memory retention
    \item Reset gate controls influence of past state
\end{itemize}

\medskip

\noindent
\textbf{Insight.}  
Gating mechanisms allow the network to dynamically control information flow.

\section{Long-Term Dependencies}

A central challenge in sequence modeling is capturing dependencies across long time spans.

\medskip

\noindent
\textbf{Problem.}
\begin{itemize}
    \item Information from early time steps may be lost
    \item Gradients may vanish during backpropagation through time
\end{itemize}

\medskip

\noindent
\textbf{Backpropagation through time (BPTT).}  
Training RNNs involves propagating gradients through many time steps:
\[
\frac{\partial L}{\partial h_t} = \frac{\partial L}{\partial h_T} \prod_{k=t}^{T-1} \frac{\partial h_{k+1}}{\partial h_k}.
\]

\medskip

\noindent
\textbf{Gradient behavior.}
\begin{itemize}
    \item Repeated multiplication can lead to vanishing or exploding gradients
\end{itemize}

\medskip

\noindent
\textbf{Role of gating.}
\begin{itemize}
    \item Provides paths for gradients to flow more easily
    \item Allows selective memory retention
\end{itemize}

\medskip

\noindent
\textbf{Practical techniques.}
\begin{itemize}
    \item Gradient clipping
    \item Truncated backpropagation
\end{itemize}

\medskip

\noindent
\textbf{Insight.}  
Modeling long-term dependencies requires both architectural and optimization considerations.

\section{Limitations of Recurrence}

Despite their strengths, recurrent architectures have important limitations.

\medskip

\noindent
\textbf{Sequential computation.}
\begin{itemize}
    \item Processing must occur step by step
    \item Limits parallelization
\end{itemize}

\medskip

\noindent
\textbf{Difficulty with long-range dependencies.}
\begin{itemize}
    \item Even gated architectures have limits
    \item Information may degrade over long sequences
\end{itemize}

\medskip

\noindent
\textbf{Optimization challenges.}
\begin{itemize}
    \item Training can be unstable
    \item Sensitive to hyperparameters
\end{itemize}

\medskip

\noindent
\textbf{Memory bottleneck.}  
A fixed-size hidden state must summarize the entire past.

\medskip

\noindent
\textbf{Insight.}  
Recurrence provides a natural framework for sequences, but may not scale effectively to complex long-range dependencies.

\section{A Unifying Perspective}

Sequence models extend neural networks to temporal data:

\begin{itemize}
    \item \textbf{State:} hidden representation evolves over time
    \item \textbf{Dynamics:} recursive updates define temporal behavior
    \item \textbf{Memory:} information is stored in hidden states
\end{itemize}

\medskip

\noindent
\textbf{Architectural evolution.}
\begin{itemize}
    \item RNN $\rightarrow$ LSTM/GRU $\rightarrow$ attention-based models
\end{itemize}

\medskip

\noindent
\textbf{Key insight.}  
The design of sequence models revolves around managing information flow across time.

\section{Outlook}

This chapter introduced sequence models:
\begin{itemize}
    \item recurrent neural networks,
    \item gated architectures,
    \item challenges of long-term dependencies,
    \item limitations of recurrence.
\end{itemize}

\medskip

In the next chapter, we study attention mechanisms, which address many of the limitations of recurrent models.

\medskip

\noindent
\textbf{Guiding principle.}  
Effective sequence modeling requires mechanisms for preserving and accessing relevant information over time.